\section{The superposition principle}

\SourcePage{16}{19}
The radical change in physical ideas about motion that quantum mechanics
requires in comparison with classical mechanics naturally calls for an equally
radical change in the theory's mathematical formalism. In this connection the
first question that arises concerns the means of describing a state in quantum
mechanics.

Let $q$ stand for the totality of coordinates of a quantum
system, and $dq$ for the product of their differentials (it
is called a volume element of the system's \textit{configuration space}); for
one particle, $dq$ coincides with the volume element $dV$ of ordinary space.

The mathematical formalism of quantum mechanics rests on the assertion that a
system's state may be described by a certain (generally complex) function of
the coordinates $\Psi(q)$, and that the squared modulus of this function
determines the probability distribution of the coordinate values:
$|\Psi|^2\,dq$ is%
\SourcePage{17}{20}
the probability that a measurement made on the system will find coordinate
values in the element $dq$ of configuration space. The function $\Psi$ is
called the system's \textit{wave function}\footnote{It was first introduced
into quantum mechanics by \textit{Schrödinger} (\textit{E.~Schrödinger},
1926).}.

Knowledge of the wave function makes it possible in principle to calculate
the probabilities of the various outcomes of any measurement in general (not
necessarily a coordinate measurement). Every one of these probabilities is
given by an expression bilinear in $\Psi$ and $\Psi^*$. In full generality,
such an expression has the form
\begin{equation}
  \iint \Psi(q)\Psi^*(q')\varphi(q,q')\,dq\,dq',
  \tag{2.1}
\end{equation}
where the function $\varphi(q,q')$ depends on the kind and outcome of the
measurement, and the integrations extend over the entire configuration space.
The probability $\Psi\Psi^*$ itself for the different coordinate values is
also an expression of this type\footnote{It is obtained from (2.1) by taking
$\varphi(q,q')=\delta(q-q_0)\delta(q'-q_0)$, where $\delta$ denotes the
so-called $\delta$-function defined below, in \secref{5}; $q_0$ denotes the
coordinate value whose probability is sought.}.

As time passes, the system's state, and with it the wave function, generally
changes. In this sense the wave function may also be regarded as a function of
time. If the wave function is known at some initial instant, then, by the very
meaning of a complete description of a state, it is thereby determined in
principle at every future instant. The actual time dependence of the wave
function is specified by equations to be derived below.

By definition, the sum of the probabilities of all possible values of the
system's coordinates must equal unity. The integral of $|\Psi|^2$ over the
whole configuration space must therefore equal unity:
\begin{equation}
  \int |\Psi|^2\,dq=1.
  \tag{2.2}
\end{equation}
This equality constitutes what is known as the \textit{normalization condition}
on wave functions. When the integral of $|\Psi|^2$ is convergent, by choosing a
suitable constant coefficient, the function $\Psi$ can always be normalized,
as it is said. We shall see below, however, that the integral of $|\Psi|^2$
may diverge, in which case $\Psi$ cannot be normalized by condition (2.2).
\SourcePage{18}{21}
Here $|\Psi|^2$ naturally does not yield the absolute coordinate
probabilities; the ratio of $|\Psi|^2$ between two distinct points of
configuration space, however, does determine the relative probability of
the coordinate values.

Since all quantities with direct physical meaning that are calculated using
the wave function have the form (2.1), in which $\Psi$ occurs multiplied by
$\Psi^*$, it is clear that a normalized wave function is defined only up to a
constant \textit{phase factor} of the form $e^{i\alpha}$, where $\alpha$ is
any real number. This ambiguity is fundamental and cannot be removed; it is,
however, immaterial, since it has no effect on any physical result.

The positive content of quantum mechanics is founded on a number of assertions
about the properties of the wave function, as follows.

Suppose that in the state with wave function $\Psi_1(q)$ a certain measurement
leads with certainty to a definite outcome, outcome 1, while in the state
$\Psi_2(q)$ it leads to outcome 2. It is then postulated that every linear
combination of $\Psi_1$ and $\Psi_2$, that is, every function of the form
$c_1\Psi_1+c_2\Psi_2$ ($c_1$, $c_2$ being constants), describes a state in
which the same measurement yields either outcome 1 or outcome 2. Moreover, if
the time dependence of the states is known, being given in one case by the
function $\Psi_1(q,t)$ and in the other by $\Psi_2(q,t)$, then any linear
combination of them also gives a possible time dependence of a state.

These assertions constitute the so-called \textit{superposition principle for
states}, the fundamental positive principle of quantum mechanics. It follows
from this principle, in particular, that every equation satisfied by wave
functions must be linear in $\Psi$.

Consider a system made up of two parts, and suppose its state is specified in
such a way that each part is described completely\footnote{This, of course,
also gives a complete description of the state of the system as a whole. We
emphasize, however, that the converse statement is by no means valid: in
general, a complete description of the state of the system as a whole does not
yet describe the states of its separate parts completely (see also \secref{14}).}.
One can then assert that the coordinate probabilities for $q_1$ in the
first part are independent of those for the coordinates $q_2$
of the second, and therefore that the probability distribution for the whole
system must equal the product of the probability distributions for its parts.
It follows that the system's wave function $\Psi_{12}(q_1,q_2)$ may be written
as a product of the wave functions $\Psi_1(q_1)$ and
$\Psi_2(q_2)$ of its%
\SourcePage{19}{22}
parts:
\begin{equation}
  \Psi_{12}(q_1,q_2)=\Psi_1(q_1)\Psi_2(q_2).
  \tag{2.3}
\end{equation}
If the two parts do not interact with each other, this relation between the
wave functions of the system and of its parts is preserved at future instants
as well:
\begin{equation}
  \Psi_{12}(q_1,q_2,t)=\Psi_1(q_1,t)\Psi_2(q_2,t).
  \tag{2.4}
\end{equation}
